\documentclass[aspectratio=169,11pt]{beamer}
% ===== 高等数学辅导课件 通用样式 =====
\usepackage[UTF8]{ctex}
\usepackage{amsmath,amssymb,amsthm}
\usepackage{fontspec}
\setCJKmainfont{Microsoft YaHei}
\setCJKsansfont{Microsoft YaHei}
\setmainfont{Microsoft YaHei}
\setsansfont{Microsoft YaHei}

% 颜色
\definecolor{navy}{RGB}{31,59,99}
\definecolor{blacktext}{RGB}{26,26,26}
\definecolor{redtext}{RGB}{192,0,0}
\definecolor{graytext}{RGB}{89,89,89}

% 去除所有模板装饰，纯白底纯内容
\setbeamertemplate{navigation symbols}{}
\setbeamertemplate{headline}{}
\setbeamertemplate{footline}{}
\setbeamertemplate{frametitle}[default][left]
\setbeamercolor{background canvas}{bg=white}
\setbeamercolor{normal text}{fg=blacktext}
\setbeamercolor{frametitle}{fg=navy}
\setbeamerfont{frametitle}{series=\bfseries,size=\Large}
\setbeamersize{text margin left=1.3cm, text margin right=1.3cm}

% 段落与行距
\linespread{1.25}
\setlength{\parskip}{0.4em}

% 自定义命令
\newcommand{\sub}[1]{{\color{navy}\bfseries #1}}
\newcommand{\con}[1]{{\color{redtext}\bfseries #1}}
\newcommand{\hint}[1]{{\color{graytext}\small #1}}
\newcommand{\blk}[1]{{\color{blacktext} #1}}


\title{高等数学辅导 · 第4课时}
\subtitle{极限运算法则、两个重要极限与等价无穷小}
\author{学生：熊张羽 \quad 授课：王老师 \quad 45分钟}
\date{}

\begin{document}

\begin{frame}
\vspace{1.2cm}
{\color{navy}\bfseries\Huge 高等数学辅导 · 第4课时}\\[1.2em]
{\color{blacktext}\bfseries\LARGE 极限运算法则、两个重要极限与等价无穷小}\\[1.6em]
{\color{graytext}
本课时内容：\\[0.3em]
知识点梳理：极限四则运算法则、复合函数极限；两个重要极限；等价无穷小替换\\
五类常见题型与解法\\
中档题精讲 4 道\\
课后习题 9 道
}
\vspace{1.4em}

{\color{graytext} 学生：熊张羽 \quad 授课教师：王老师 \quad 45分钟}
\end{frame}

\begin{frame}{一、知识点梳理}
\sub{1. 极限四则运算法则}
若 $\lim f=A,\ \lim g=B$，则
\[
\lim(f\pm g)=A\pm B,\quad \lim(fg)=AB,\quad \lim\frac fg=\frac AB\ (B\ne 0).
\]
推论：$\lim[cf(x)]=c\lim f(x)$，$\lim[f(x)]^n=[\lim f(x)]^n$。

\sub{2. 复合函数的极限}
若 $\lim_{x\to x_0}\varphi(x)=u_0$，且 $\lim_{u\to u_0}f(u)=A$，则
\[
\lim_{x\to x_0}f(\varphi(x))=A.
\]
即极限符号与函数符号可交换：$\lim f(\varphi(x))=f(\lim\varphi(x))$（要求 $f$ 连续）。

\sub{3. 两个重要极限}
\[
\lim_{x\to 0}\frac{\sin x}{x}=1,
\qquad
\lim_{x\to\infty}\left(1+\frac1x\right)^x=e
\quad\left(\text{或}\ \lim_{x\to 0}(1+x)^{1/x}=e\right).
\]
\end{frame}

\begin{frame}{一、知识点梳理（续）}
\sub{4. 无穷小的比较}
设 $\alpha,\beta$ 为同一过程中的无穷小。
\[
\lim\frac\beta\alpha=
\begin{cases}
0, & \beta\text{ 是 }\alpha\text{ 的高阶无穷小，记 }\beta=o(\alpha),\\
c\ne 0, & \beta\text{ 与 }\alpha\text{ 同阶无穷小},\\
1, & \beta\text{ 与 }\alpha\text{ 等价无穷小，记 }\beta\sim\alpha,\\
\infty, & \beta\text{ 是 }\alpha\text{ 的低阶无穷小}.
\end{cases}
\]

\sub{5. 常用等价无穷小（$x\to 0$）}
\[
\sin x\sim x,\ \tan x\sim x,\ \arcsin x\sim x,\ \arctan x\sim x,
\]
\[
1-\cos x\sim \frac{x^2}{2},\ e^x-1\sim x,\ \ln(1+x)\sim x,\ (1+x)^\alpha-1\sim \alpha x.
\]

\sub{6. 等价无穷小替换定理}
若 $\alpha\sim\alpha',\ \beta\sim\beta'$，则 $\lim\dfrac{\beta f}{\alpha}=\lim\dfrac{\beta' f}{\alpha'}$。
\hint{关键：只能替换乘除因子，不能替换加减项！}
\end{frame}

\begin{frame}{二、常见题型与解法}
\sub{题型1 \quad $0/0$ 型极限}
解法：因式分解约零因子 / 根式有理化 / 等价无穷小替换 / 第一个重要极限。

\sub{题型2 \quad $\infty/\infty$ 型极限（$x\to\infty$）}
解法：分子分母同除以最高次幂，利用"无穷大倒数为无穷小"。

\sub{题型3 \quad $\infty-\infty$ 型极限}
解法：能通分则通分；含根式之差则乘以共轭根式有理化。

\sub{题型4 \quad $1^\infty$ 型极限}
解法：凑成 $(1+\square)^{1/\square}$ 形式用第二个重要极限，或取对数。

\sub{题型5 \quad 等价无穷小替换}
解法：在乘除结构中，将复杂无穷小替换为简单等价无穷小。注意加减法不能直接替换。
\end{frame}

\begin{frame}{三、中档题 · 例1}
\sub{例1} 求 $\displaystyle\lim_{x\to 0}\frac{\tan x-\sin x}{x^3}$。

\blk{解：} 先恒等变形（加减法中不能直接替换 $\tan x\sim x,\ \sin x\sim x$，否则得 $0$，错误）：
\[
\tan x-\sin x=\sin x\left(\frac1{\cos x}-1\right)=\frac{\sin x(1-\cos x)}{\cos x}.
\]
当 $x\to 0$ 时 $\sin x\sim x,\ 1-\cos x\sim \dfrac{x^2}{2},\ \cos x\to 1$，故
\[
\lim_{x\to 0}\frac{\tan x-\sin x}{x^3}
=\lim_{x\to 0}\frac{x\cdot \frac{x^2}{2}}{x^3\cdot \cos x}
=\frac12.
\]

\con{结论：原式 $=\dfrac12$。}

\hint{警示：等价无穷小替换只用于乘除因子，不用于加减项。}
\end{frame}

\begin{frame}{三、中档题 · 例2}
\sub{例2} 求 $\displaystyle\lim_{x\to\infty}\left(\frac{2x-1}{2x+1}\right)^x$。

\blk{解：} 这是 $1^\infty$ 型。将底数化为 $1+\square$ 的形式：
\[
\frac{2x-1}{2x+1}=1+\frac{-2}{2x+1}.
\]
利用第二个重要极限的推广 $\lim(1+u)^{1/u}=e\ (u\to 0)$：
\[
\left(\frac{2x-1}{2x+1}\right)^x
=\left[\left(1+\frac{-2}{2x+1}\right)^{\frac{2x+1}{-2}}\right]^{\frac{-2x}{2x+1}}.
\]
方括号内 $\to e$，指数 $\dfrac{-2x}{2x+1}\to -1$，故
\[
\lim_{x\to\infty}\left(\frac{2x-1}{2x+1}\right)^x=e^{-1}.
\]

\con{结论：原式 $=e^{-1}$。}
\end{frame}

\begin{frame}{三、中档题 · 例3}
\sub{例3} 求 $\displaystyle\lim_{x\to 0}\frac{\ln(1+2x)}{\sin 3x}$。

\blk{解：} 当 $x\to 0$ 时，$\ln(1+2x)\sim 2x$，$\sin 3x\sim 3x$。

二者均为乘除因子，可直接替换：
\[
\lim_{x\to 0}\frac{\ln(1+2x)}{\sin 3x}
=\lim_{x\to 0}\frac{2x}{3x}
=\frac23.
\]

\con{结论：原式 $=\dfrac23$。}

\hint{注意：$\ln(1+2x)\sim 2x$ 是由 $\ln(1+u)\sim u$ 令 $u=2x$ 得到的；$\sin 3x\sim 3x$ 同理。}
\end{frame}

\begin{frame}{三、中档题 · 例4}
\sub{例4} 求 $\displaystyle\lim_{x\to+\infty}\left(\sqrt{x^2+x}-\sqrt{x^2-x}\right)$。

\blk{解：} 这是 $\infty-\infty$ 型，含根式之差，乘以共轭根式有理化：
\[
\sqrt{x^2+x}-\sqrt{x^2-x}
=\frac{(\sqrt{x^2+x}-\sqrt{x^2-x})(\sqrt{x^2+x}+\sqrt{x^2-x})}{\sqrt{x^2+x}+\sqrt{x^2-x}}
\]
\[
=\frac{(x^2+x)-(x^2-x)}{\sqrt{x^2+x}+\sqrt{x^2-x}}
=\frac{2x}{\sqrt{x^2+x}+\sqrt{x^2-x}}.
\]
分子分母同除以 $x$（注意 $x>0$）：
\[
=\frac{2}{\sqrt{1+\frac1x}+\sqrt{1-\frac1x}}
\to \frac{2}{1+1}=1.
\]

\con{结论：原式 $=1$。}
\end{frame}

\begin{frame}{四、课后习题（1—5）}
\begin{enumerate}
  \item 求 $\displaystyle\lim_{x\to 1}\frac{x^2-1}{2x^2-x-1}$。
  \item 求 $\displaystyle\lim_{x\to 0}\frac{\sin 3x}{\sin 5x}$。
  \item 求 $\displaystyle\lim_{x\to 0}\frac{1-\cos x}{x^2}$。
  \item 求 $\displaystyle\lim_{x\to 0}\frac{e^{2x}-1}{\ln(1+x)}$。
  \item 求 $\displaystyle\lim_{x\to 0}\frac{\arctan x}{x}$。
\end{enumerate}
\end{frame}

\begin{frame}{四、课后习题（6—9）}
\begin{enumerate}\setcounter{enumi}{5}
  \item 求 $\displaystyle\lim_{x\to 0}(1+3x)^{2/x}$。
  \item 求 $\displaystyle\lim_{x\to 0}\frac{\sqrt[3]{1+x^2}-1}{x^2}$。
  \item 求 $\displaystyle\lim_{x\to+\infty}\left(\sqrt{x^2+x}-\sqrt{x^2-x}\right)$。
  \item 求 $\displaystyle\lim_{x\to 0}\frac{\tan x-\sin x}{\sin^3 x}$。
\end{enumerate}
\end{frame}

\begin{frame}{课后任务}
\begin{itemize}
  \item 完成本课时课后习题 1—9。
  \item 重点掌握：等价无穷小替换（仅限乘除）、$1^\infty$ 型凑第二个重要极限、$\infty-\infty$ 有理化。
  \item 下节课内容：极限计算综合与错题强化（摸底 Q13、Q14 同类题精讲）。
\end{itemize}
\end{frame}

\end{document}
